\subsection{Geographic extent of simultaneous humid heat}

The attenuation of broad spatial contrasts accompanies a larger connected
humid-heat footprint (Figure~\ref{fig:application-coverage}). At
24$^\circ$C, the high-minus-middle differences in total and largest connected
area are 12.91 and 12.57 percentage points, respectively, each corresponding
to approximately 400,000 km$^2$ of represented land. The coverage contrast
therefore includes substantial differences in the extent of contiguous
humid-heat regions.

Threshold dependence is central to this interpretation. The area separation
is greatest around 23--24$^\circ$C and narrows towards the highest thresholds.
The geographic expression of high regional humid heat consequently depends
on the temperature level being considered. Spatial contraction alone does
not determine this curve: regional mean and spatial pattern change together
between the groups. Direct coverage measurement establishes the extent of
their combined association.

\begin{figure}[H]
\centering
\includegraphics[width=\textwidth]{application_coverage.pdf}
\caption{High-minus-middle differences in simultaneous humid-heat coverage
at the original daily peak times. (a) Total exceedance area. (b) Largest
four-neighbour connected exceedance area. Thick curves average months and
33 evaluation summers equally; thin grey curves show every summer. Shading
gives pointwise 95\% percentile intervals from 1,999 resamples of five-year
calendar blocks. All 77 development-defined temperature thresholds are shown.
Areas refer to the 3.221-million-km$^2$ represented land support.}
\figalttext{Both area-difference curves rise from small positive values at low
thresholds to approximately thirteen percentage points near 24 degrees
Celsius, then decline toward the upper end of the threshold range. Grey
annual curves show variation between summers, including some negative values
at high thresholds.}
\label{fig:application-coverage}
\end{figure}

The comparison of spatial summaries identifies where most coverage
information resides (Table~\ref{tab:application-comparison}). The principal
reduction in error occurs when the regional mean is supplemented by simple
spatial summaries. Distance-resolved predictors provide small further
improvements, with graph and variogram representations performing similarly.
The same conclusion holds for the additional-site reference. Thus coverage
estimation and structural interpretation have different informational
requirements: the simple summaries describe geographic extent well, while
the graph profile connects regional dispersion to its spatial scales and
components. The magnitude and annual consistency of the incremental gains
are documented in Supplementary Section~S13.

\begin{table}[!t]
\centering\small
\caption{Held-out mean absolute error in area percentage points. Errors are
averaged over all 77 thresholds (9.5--28.5$^\circ$C at 0.25$^\circ$C steps),
then equally over months and 33 evaluation summers. All-day errors use 3,036
June--August daily fields in 1991--2025 excluding 2015 and 2022; high-day errors
use the 792 original upper-quartile days. The last two columns use the 344
additional sites. All inputs come from the original 121 sites.}
\label{tab:application-comparison}
\begin{tabular}{@{}lrrrrrr@{}}
\toprule
& \multicolumn{2}{c}{Exceedance area} & \multicolumn{2}{c}{Largest component} & \multicolumn{2}{c}{Additional sites}\\
\cmidrule(lr){2-3}\cmidrule(lr){4-5}\cmidrule(lr){6-7}
Method & All & High & All & High & All & High\\
\midrule
Mean and season & 2.794 & 2.291 & 3.070 & 2.593 & 2.818 & 2.314 \\
Strong simple baseline & 0.827 & 0.747 & 1.389 & 1.268 & 1.028 & 0.929 \\
Simple + variogram & 0.815 & 0.741 & 1.366 & 1.266 & 1.013 & 0.922 \\
Simple + graph & 0.820 & 0.743 & 1.379 & 1.265 & 1.020 & 0.925 \\
Spatial interpolation & 1.398 & 1.414 & 1.823 & 1.859 & 1.861 & 1.882 \\
\bottomrule
\end{tabular}
\end{table}

